% =====================================================================
%  1bat-ud1-11.tex  —  HORA 11: Per què la notació científica? (descoberta)
% =====================================================================
\horatitol{Sessió 11 — Per què necessitem la notació científica?}%
{Descobrir-ho amb pressupostos reals: una manera còmoda d'escriure i comparar xifres molt grans.}

\subsection{Idea clau}

Canviem de bloc. Fins ara hem treballat llindars i edats; ara mirem
\textbf{xifres molt grans} de l'àmbit social: pressupostos públics, poblacions. Per
exemple, el pressupost del Departament de Drets Socials o el nombre de persones en
risc d'exclusió. Escriure tants zeros és incòmode i fàcil d'equivocar: la
\textbf{notació científica} ho resol.

\begin{idea}
\textbf{Notació científica:} tot nombre gran s'escriu com
\[
a\times 10^{n},
\]
on $a$ és la \textbf{mantissa} (un nombre $\ge 1$ i $<10$, és a dir, amb una sola
xifra abans de la coma) i $n$ és l'\textbf{exponent} (enter).
\begin{itemize}[leftmargin=1.4em, itemsep=2pt]
  \item L'exponent $n$ \textbf{conta quantes posicions} es mou la coma.
  \item Per a nombres grans, $n$ és \textbf{positiu}: $\num{3000}=3\times 10^{3}$.
  \item A la calculadora apareix com \texttt{3E3} o \texttt{3\,EXP\,3}.
\end{itemize}
\textbf{Exemple guia:} $\num{3326000000}=3,326\times 10^{9}$ (la coma s'ha mogut $9$
llocs cap a l'esquerra).
\end{idea}

\subsection{Exemples}

\begin{exemple}[el pressupost de Drets Socials]
El pressupost és d'uns $\num{3326}$ milions d'euros, és a dir
$\num{3326000000}\eur$. Comptem quants llocs hem de moure la coma fins a deixar una
sola xifra al davant ($3,326$):
\[
\num{3326000000} = 3{,}326\times 10^{9}\eur.
\]
Molt més còmode que escriure els nou zeros, i més difícil d'equivocar.
\end{exemple}

\begin{exemple}[la població en risc d'exclusió]
D'una població d'uns $8$ milions d'habitants, un $24,4\%$ està en risc d'exclusió.
Primer el càlcul:
\[
0,244\times\num{8000000}=\num{1952000}\text{ persones}.
\]
En notació científica, movem la coma fins a deixar una xifra al davant ($1,952$):
\[
\num{1952000}\approx 1{,}95\times 10^{6}\text{ persones}.
\]
\end{exemple}

\subsection{Exercicis resolts}

\begin{exresolt}[passar a notació científica]
Escriu en notació científica: (a) $\num{45000}$; \ (b) $\num{2300000}$; \
(c) $\num{870}$.
\solucio
Movem la coma fins a deixar una xifra (entre $1$ i $9$) al davant, i comptem els llocs:
\begin{enumerate}[label=\alph*), leftmargin=1.6em]
  \item $\num{45000}=4,5\times 10^{4}$ (la coma es mou $4$ llocs).
  \item $\num{2300000}=2,3\times 10^{6}$ (6 llocs).
  \item $\num{870}=8,7\times 10^{2}$ (2 llocs).
\end{enumerate}
\resposta{(a) $4,5\times 10^{4}$; (b) $2,3\times 10^{6}$; (c) $8,7\times 10^{2}$.}
\end{exresolt}

\begin{exresolt}[de notació científica a nombre normal]
Escriu com a nombre «normal»: (a) $5\times 10^{5}$; \ (b) $1,2\times 10^{7}$.
\solucio
L'exponent diu quants llocs movem la coma cap a la \textbf{dreta} (afegint zeros):
\begin{enumerate}[label=\alph*), leftmargin=1.6em]
  \item $5\times 10^{5}=\num{500000}$ (5 zeros).
  \item $1,2\times 10^{7}=\num{12000000}$ (la coma es mou $7$ llocs).
\end{enumerate}
\resposta{(a) $\num{500000}$; (b) $\num{12000000}$.}
\end{exresolt}

\subsection{Exercicis per fer}

\begin{exercicis}
\begin{exlist}
  \item Escriu en notació científica:
        \begin{enumerate}[label=\alph*), leftmargin=1.6em, topsep=2pt]
          \item $\num{6000}$ \hfill \item $\num{52000}$ \hfill \item $\num{1500000}$
          \hfill \item $\num{340}$
        \end{enumerate}
  \item Escriu com a nombre normal:
        \begin{enumerate}[label=\alph*), leftmargin=1.6em, topsep=2pt]
          \item $7\times 10^{4}$ \hfill \item $3,5\times 10^{6}$ \hfill
          \item $9,1\times 10^{3}$
        \end{enumerate}
  \item El pressupost d'una partida social és de $\num{1280000000}\eur$. Escriu-lo en
        notació científica.
  \item Una ciutat té $\num{1600000}$ habitants. Escriu la xifra en notació
        científica i digues quantes xifres té el nombre original.
  \item \textbf{(Estimació)} Sense fer el càlcul exacte, quants zeros (aproximadament)
        tindrà el resultat de multiplicar $2\times 10^{3}$ per $4\times 10^{5}$? (Pista:
        suma els exponents.)
  \item \textbf{(Context real)} El $24,4\%$ d'una població de $8$ milions està en risc
        d'exclusió. Calcula quantes persones són i escriu el resultat en notació
        científica (arrodonit a dues xifres a la mantissa).
\end{exlist}
\end{exercicis}

% ---------------------------------------------------------------------
\fullsolucions{Sessió 11}
\begin{solucions}
\begin{sollist}
  \item (a) $6\times 10^{3}$; \quad (b) $5,2\times 10^{4}$; \quad
        (c) $1,5\times 10^{6}$; \quad (d) $3,4\times 10^{2}$.
  \item (a) $\num{70000}$; \quad (b) $\num{3500000}$; \quad (c) $\num{9100}$.
  \item $\num{1280000000}=1,28\times 10^{9}\eur$.
  \item $\num{1600000}=1,6\times 10^{6}$ habitants; el nombre original té $7$ xifres.
  \item Els exponents sumen $3+5=8$, i $2\times 4=8$, així que el resultat és
        $8\times 10^{8}$: aproximadament $8$ seguit de $8$ zeros.
  \item $0,244\times\num{8000000}=\num{1952000}\approx 1,95\times 10^{6}$ persones.
\end{sollist}
\end{solucions}
